\documentclass[12pt,a4paper,oneside]{extarticle}
% --------------------------------------------------------
% for PDF/A-1b pdf file - +++ DO NOT REMOVE THIS +++
% https://en.wikipedia.org/wiki/PDF/A
\PassOptionsToPackage{hidelinks}{hyperref}
\usepackage[a-1b]{pdfx}
\usepackage{hyperref}
\hypersetup{
	colorlinks=true,
	allbordercolors = white,
	%linkbordercolor = {white},
	%allcolors = black,
	citecolor = black,
	urlcolor = blue,
	linkcolor = blue,
	anchorcolor = black,
	%citecolor = green,
	filecolor = cyan,
	menucolor = red,
	%urlcolor = magenta,
	%linkcolor=blue,
	linktocpage=true
}

\usepackage{graphicx} % Required for inserting images
\usepackage{pdfpages}
\usepackage{enumerate}
%\usepackage[top=45mm,bottom=35mm,inner=35mm,outer=25mm]{geometry} # hmarginratio=1:1,
% \usepackage[top=25mm, bottom=25mm, left=25mm, right=25mm, columnsep=25pt]{geometry} % Document margins
% \usepackage[top=25mm, bottom=25mm, left=20mm, right=20mm, columnsep=25pt,headsep=1cm,headheight=50.5pt,footskip=2cm]{geometry} % Document margins
\usepackage[top=3.5cm, bottom=2.5cm, left=20mm, right=20mm, marginparsep=0.35cm,headsep=1cm,headheight=1.6cm,footskip=2cm]{geometry}
\usepackage{setspace}
% \setstretch{1.}
\usepackage{url}
\usepackage{comment}

\renewcommand{\refname}{}
\usepackage{fancyhdr}
\setlength{\headheight}{50.5pt}
% \setlength{\footskip}{16mm}

\pagestyle{fancy}

\fancyhead{}
\fancyfoot{}

\fancypagestyle{default_page}{
	\fancyhead{}
	\fancyfoot{}
	\fancyhead[L]{\includegraphics[width=0.3\textwidth]{../LOGO_UNISI_ORIZZONTALE_NERO_0.pdf}}
	\fancyhead[R]{
		\textit{Dipartimento di Scienze Fisiche, della Terra e dell’Ambiente} \\ \textbf{Dottorato in Fisica Sperimentale}
		}
	\fancyfoot[L,R]{}
	\fancyfoot[C]{Via Roma, 56 $-$ 53100 Siena (Italia)}
	\renewcommand{\headrulewidth}{0.1mm}
	\renewcommand{\footrulewidth}{0.1mm}
}

% \fancyfoot[L]{Valeria D'Amante}
% \fancyfoot[C]{E1A}


%-------Definition of \myrule--
\def\myrule#1#2#3{{\hskip#1in{\hbox to #2in%
{\leaders\hbox to .00625in{\hfil.\hfil}\hfill}}%
\par\hskip#1in#3\vskip1cm}}
%------------------------------

\begin{document}
\pagestyle{default_page}

\section*{REPORT ON THE PhD ACTIVITIES CARRIED OUT BY \textit{Valeria D'Amante}}


The doctoral student Valeria D'Amante, enrolled in the PhD programme in “Fisica Sperimentale” at the University of Siena, cycle XXXVI, carried out the following activities during the PhD course:

\subsection*{RESEARCH ACTIVITIES}
\begin{enumerate}
    \item \textbf{Double Higgs Production (CMS Experiment, CERN, Geneva)}\\
    Conducted an in-depth analysis of resonant and non-resonant Higgs boson pair production ($HH \to b\bar{b}\tau\tau$) using $ \sqrt{s} = 13\,\textrm{TeV} $ proton-proton collision data collected during CMS Run 2. This work focused on the search for resonances in the mass range of 250 GeV to 3 TeV, targeting scenarios predicted by Beyond Standard Model (BSM) theories. Developed an independent analysis framework leveraging modern tools specifically tailored for high-energy physics, as ROOT RDataFrames and the Luigi Analysis Workflow (LAW). Contributed to signal and background optimization through advanced machine learning techniques, improving event classification and sensitivity. Performed detailed evaluations of systematic uncertainties and statistical limit extractions, which led to the establishment of upper bounds on Higgs boson pair production cross-sections in both SM and BSM contexts. The results were included in the CMS combined $HH$ analysis published in \textit{Nature}, commemorating the 10th anniversary of the Higgs boson discovery.

    \item \textbf{Machine Learning for Hadronic Tau ($\tau_h$) Identification and Reconstruction (CMS Experiment, CERN)}\\
    Development of the L2TauNNTag algorithm to replace traditional cut-based methods in High-Level Trigger (HLT) paths. Demonstrated improved efficiency and purity in selecting hadronic taus, contributing to Run 3 data-taking.

    \item \textbf{Measurement of CKM Matrix Elements (CMS Experiment, CERN)}\\
    Model-independent measurement of the CKM matrix elements $|V_{tq}|$ (where $q = d, s, b$) using single top production in the $t$-channel. Implemented machine learning techniques (DNN and BDT) to enhance signal extraction and top quark identification.

    \item \textbf{Neutrino-Argon Interactions (SBND Experiment, Fermilab)}\\
    Analysis of neutrino interaction rates in the Short-Baseline Near Detector (SBND) at Fermilab. Focused on specific final-state topologies to study nuclear effects in neutrino-argon interactions.

    \item \textbf{Dark Matter Direct Detection Sensitivity (DarkSide-20k, LNGS)}\\
    Investigated WIMP interactions in liquid argon and xenon detectors, producing exclusion plots under the null-result hypothesis for experiments such as DarkSide-20k and Xenon-1T.
\end{enumerate}
% \newpage
% \vspace{1cm}
\subsection*{DOCTORAL THESIS}

\textbf{Title}: Search for resonant double Higgs production in the bb$\tau\tau$ final state at the CMS experiment.
\\
\textbf{Supervisors}: Maria Agnese Ciocci
\\
\textbf{Description of the issues and main results}:
The search for resonant double Higgs production in the $HH \to b\bar{b}\tau\tau$ final state explores the possibility of Beyond Standard Model (BSM) resonances, such as spin-0 or spin-2 particles, decaying into Higgs boson pairs. These resonances are searched for in the mass range of 250 GeV to 3 TeV, using $ \sqrt{s} = 13\,\textrm{TeV} $ proton-proton collision data from CMS Run 2.\\
\textit{Challenges:}
\begin{itemize}
    \item Significant background contributions from $t\bar{t}$, $Z + \textrm{jets}$, and single Higgs production.
    \item Limited statistics due to the low $HH$ cross-section.
    \item Accurate reconstruction of hadronically decaying $\tau$-leptons ($\tau_h$).
    \item Systematic uncertainties from detector performance and background modeling.
\end{itemize}
\noindent \textit{Results:}
Optimized selection criteria, including machine learning-based event classifiers, were employed to enhance signal-background separation. A shape-based statistical analysis was performed, placing upper limits on the resonance production cross-section as a function of its mass. Constraints were set on BSM models predicting heavy resonances decaying to $HH$, improving or complementing existing limits in the explored range. A novel analysis framework was developed, enabling efficient data processing and validation.\\ %  incorporating CMS-specific tools (e.g., NanoAOD, ROOT RDataFrames, and Luigi workflows)
\textit{Impact:}
Although no significant deviations from Standard Model predictions were observed, the results contribute to ongoing BSM searches and pave the way for improved sensitivity in future LHC runs. The advanced techniques for $\tau_h$ reconstruction and event classification developed in this work are essential for probing the Higgs sector with increased data.

% \newpage
\subsection*{PARTICIPATION IN CONGRESSES, SCHOOLS AND COURSES}

\subsubsection*{Conferences}
\begin{itemize}
    \item \textbf{GravityShapePisa (GRASP) 2023}, Pisa, Italy, October 24–27, 2022.
    \item \textbf{CMS Italia National Meeting}, Turin, Italy, October 18–20, 2023.
    \item \textbf{Higgs 2022}, Pisa, Italy, November 6–11, 2022.
    \item \textbf{Higgs@CMS Italia Workshop}, CERN, Geneva, May 11–12, 2022.
    \item \textbf{2022 Higgs Workshop}, CERN, Geneva, March 28–30, 2022.
    \item \textbf{Higgs 2021}, Online Conference, October 18–22, 2021.
    \item \textbf{PyHEP Workshop}, CERN, Online, July 5–9, 2021.
    \item \textbf{11th CMS Induction Course}, CERN, Geneva, March 10–12, 2021.
\end{itemize}


\subsubsection*{Talks and International Workshops}
\begin{itemize}
    \item \textbf{The 17th International Workshop on Tau Lepton Physics (TAU2023)}, Louisville, USA, December 4–8, 2023.\\
    Talk: "Searches for New Physics that Couple with Third Generation Fermions."
    \item \textbf{HH2022: Higgs Hunting 2022}, IJCLab, Paris, September 12–14, 2022.\\
    Presentation: "Search for Non-Resonant Higgs Boson Pair Production in the Final State with Two Bottom Quarks and Two Tau Leptons."
    \item \textbf{ICHEP2022: 41st International Conference on High Energy Physics}, Bologna, Italy, July 6–13, 2022.\\
    Talk: "Search for Resonant and Nonresonant Di-Higgs Boson Production at CMS Using Jet Substructure Techniques."
    \item \textbf{CMS Week (Tau POG)}, CERN, Geneva, September 19–23, 2022.\\
    Presentation: "Tau Trigger Summary."
    \item \textbf{INFIERI Summer School}, Madrid, August 22–September 5, 2021.\\
    Poster: "A Machine Learning Algorithm for Tau Leptons Identification at L2 Trigger in the CMS Experiment."
\end{itemize}

\subsubsection*{International Schools}
\begin{itemize}
    \item \textbf{ESHEP 2023: European School of High-Energy Physics}, Grenaa, Denmark, September 6–19, 2023.
    \item \textbf{INFN School of Statistics 2022}, Paestum, Italy, May 15–20, 2022.
    \item \textbf{INFIERI Summer School}, Madrid, August 22–September 5, 2021.
    \item \textbf{XXIV CMS Data Analysis School}, Fermilab, January 5–16, 2021.
    \item \textbf{"Re-writing Nuclear Physics Textbooks" Summer School}, University of Pisa, July 19–25, 2015.
\end{itemize}

\subsection*{PUBLICATIONS}
% During the PhD period, Dr. __________ produced the following articles:
Since February 2021, the inclusion in the list of authors of the CMS Collaboration has resulted in authorship of 232 published papers, with an h-index of 46, as documented on \href{https://inspirehep.net/authors/1821915}{Inspire}.

The following list includes selected publications and papers (including CMS's ``Physics Analysis Summaries'' and ``Detector Performance Notes'') relevant:
\vspace{-1.5cm}
\bibliographystyle{abbrv}
\bibliography{biblio}
\nocite{bbttpaper}
\nocite{CMS-PAS-HIG-20-010}
\nocite{DAmante:2022xmb}
\nocite{CMS-DP-2023-024}

Additionally, a proceeding for the 17th International Workshop on Tau Lepton Physics (TAU2023), held from December 3 to 8, 2023, has been submitted and will be published online through SciPost upon assignment of a referee.\\

% \vspace{1cm}
% Date: 27/12/2024
% \vspace{2cm}

\begin{tabular}{p{3cm} p{7cm} p{3cm}}
    \textbf{Date:} &\hspace{-2cm} 27/12/2024 & \\
    &  &\\
    &  &\\
    & \hspace{1cm} \textbf{PhD student's signature} &\vspace{-6mm} \hspace{1cm} \underline{\hspace{5cm}} \\
    & &\\
    & &\\
    & &\\
    & \hspace{4.5cm} \textbf{Visto} & \vspace{-6mm} \hspace{1cm} \underline{\hspace{5cm}} \\
    & \hspace{4.1cm} \vspace{-1cm}(Il Tutore)  &\\
\end{tabular}

% \myrule{4}{2}{PhD student's signature}

% \vspace{2cm}
% \myrule{4}{2}{Visto}
% \begin{flushright}(Il tutore)\end{flushright}



% 			PhD student's signature	___________________________________

% 					Visto		___________________________________
% 					(Il Tutore)




\end{document}